\documentclass[10pt,a4paper]{amsart}
\usepackage[margin=19mm]{geometry}
\usepackage{amsmath,amssymb,mathtools}
\usepackage[scr=rsfs,cal=euler]{mathalfa}
\usepackage{xfrac}
\usepackage[unicode,hidelinks,bookmarks=false]{hyperref}
\newcommand{\ie}{i.e.\ }
\newcommand{\cf}{cf.\ }
\newcommand{\resp}{respectively\ }
\newcommand{\Subsection}{\subsection}
\providecommand{\nobreakdash}{\nobreak\hbox{-}}
\setlength{\emergencystretch}{2em}
\multlinegap0pt
\usepackage{amssymb}
\usepackage[shortlabels]{enumitem}
\usepackage{mathtools}
\newtheorem{proposition}{Proposition}
\newcommand{\depth}{\operatorname{depth}}
\newcommand{\edim}{\operatorname{edim}}
\newcommand{\pdim}{\operatorname{proj\,dim}}
\title{\Large Commutative algebra inspired by modularity lifting}
\author{Srikanth B. Iyengar}
\date{Source: arXiv:2510.11875v2 (CC BY 4.0)}
\begin{document}
\maketitle

\noindent\emph{Source and licence.} \Large Commutative algebra inspired by modularity lifting. Version arXiv:2510.11875v2 (CC BY 4.0), distributed by arXiv under the Creative Commons Attribution 4.0 International licence, http://creativecommons.org/licenses/by/4.0/, as recorded in the arXiv licence metadata for this exact version and shown on the abstract page https://arxiv.org/abs/2510.11875v2. Full licence notice: \texttt{licenses/arxiv-2510.11875-CC-BY-4.0.txt}.\par\smallskip

\noindent\textit{Editorial scope.} Counted unit: the article's proposition pr:diamond (source0.tex lines 174--177). The article's complete original proof of it is delivered with it (source0.tex lines 179--185, one \texttt{proof} environment of 7 lines). No mathematics is written, repaired or re-derived: the statement and the proof are the article's own lines, in the article's own notation, and no retained line is substituted or re-flowed.  No further source-local context is needed: every cross-reference of every retained line resolves inside the two delivered spans.  Nothing was omitted from any retained span, so the package carries no omission record.  The retained lines cite 1 works, whose entries are transcribed verbatim from the article's own bibliography source in the e-print and delivered with short editorial labels recorded in the item's reference mapping.  No retained line contains a figure environment, an image-inclusion command or drawing code, so no image is delivered and the package declares no asset dependency.  Editorial formatting only, in addition: an AMS \texttt{amsart} preamble that restates the article's own declarations and macros used by the retained lines, namely the environments \texttt{proposition} on separate counters, together with the standard packages those lines need.  The retained environments print in this document as Proposition 1 in order of appearance, whereas the article prints the same items with the numbering of its own full document.  The complete native source of the article as distributed in the arXiv e-print for this exact version is bundled unchanged as \texttt{sources/186998/source0.tex}.
\begin{proposition}
\label{pr:diamond}
Let $\varphi\colon A\to B$ be a local homomorphism of noetherian local rings, with $A$ regular. Let $N$ be a nonzero $B$-module that is finitely generated as an $A$-module and satisfies $\pdim_AN\le \edim A-\edim B$, then $N$ is free as $B$-module and the local ring $B$ is regular as well.
\end{proposition}

\begin{proof}
By the Auslander-Buchsbaum formula~\cite[Theorem~1.3.3]{Bruns/Herzog:1998a} the hypotheses on $N$ implies the first inequality below
\[
\edim B \le \depth_A N \le \depth_B N \le \dim B\le\edim B\,.
\]
The other inequalities are standard; see ~\cite[\S1.2]{Bruns/Herzog:1998a}. Thus we conclude that $\edim B=\dim B$ and hence that $B$ is regular; see \cite[\S2.2]{Bruns/Herzog:1998a}. Moreover we get also that $\depth_BN=\dim B$ which, given that $B$ is regular, implies that $N$ is free as a $B$-module, again by the Auslander-Buchsbaum formula.
\end{proof}

\begin{thebibliography}{99}
\bibitem[Bruns/]{Bruns/Herzog:1998a} {\sc W.~Bruns and J.~Herzog}, {\em {Cohen--Macaulay rings}}, vol.~39 of Cambridge Studies in Advanced Mathematics, Cambridge University Press, 1998. \newblock Revised edition.
\end{thebibliography}
\end{document}
